\section{Financial data structures: bars, multi-product series and event sampling}

\begin{qfmethod}[Activity-sampled bars: time, tick, volume and dollar bars, what each corrects, and what the evidence on their returns shows]\label{meth:activity-sampled-bars}
\qffield{Statement} \emph{Bars} are the rows of a table into which unstructured trade data are regularised; standard bar methods ``transform a series of observations that arrive at irregular frequency (often referred to as `inhomogeneous series') into a homogeneous series derived from regular sampling'' \citep[chap.~2, sec.~2.3.1, p.~26]{alpha-afml}. Each bar records a timestamp, VWAP, open, close, high, low and volume. \emph{Time bars} sample at fixed clock intervals; \emph{tick bars} each time a fixed number of transactions has taken place (for example 1,000); \emph{volume bars} each time a fixed number of the security's units has traded; \emph{dollar bars} each time a fixed market value has traded, in the security's currency (secs.~2.3.1.1--2.3.1.4, pp.~26--28).

\qffield{Conditions} López de Prado's case against time bars is that markets do not process information at a constant rate, so time bars ``oversample information during low-activity periods and undersample information during high-activity periods'', and that time-sampled series ``often exhibit poor statistical properties, like serial correlation, heteroscedasticity, and non-normality of returns'' (p.~26, citing Easley, López de Prado and O'Hara 2012). Tick bars synchronise sampling with ``a proxy of information arrival'' (p.~27), but auction prints and order fragmentation distort the tick count: an opening or closing auction is published as one trade ``for an outsized amount'' that ``could be the equivalent of thousands of ticks'', and one buy of 10 lots against ten resting orders of 1 lot prints as ten ticks; matching engines may also ``split one fill into multiple artificial partial fills'' (p.~27).

Volume bars avoid the fragmentation problem. Dollar bars keep a constant traded value when prices move a lot: after a 100 percent rise, \$1{,}000 buys half the shares (p.~28). After split adjustment, issuance and buybacks still change tick and volume counts, and dollar bars ``tend to be robust'' to those actions; the bar size can be set dynamically from free-float capitalisation or outstanding debt (p.~28).

\qffield{Practice} The evidence cited is of different kinds and does not settle which clock is best. López de Prado quotes Mandelbrot and Taylor (1967), ``Price changes over a fixed number of transactions \emph{may} have a Gaussian distribution'', says that ``multiple studies have confirmed'' that sampling by trading activity gives ``returns closer to IID Normal'' (citing Ané and Geman 2000), and reports Clark (1973) as finding volume sampling closer to IID Gaussian than tick sampling (p.~27). Ané and Geman's abstract claims that normality ``can be recovered through a stochastic time change'' and that ``the cumulative number of trades is a better stochastic clock than the volume for generating virtually perfect normality in returns'' \citep[abstract]{ane-geman-2000}, which ranks trades above volume, the reverse of the ranking reported for Clark.

\emph{A replication critique}. Murphy and Izzeldin's Monte Carlo experiments show that Ané and Geman's procedure cannot accurately recover the third and higher moments of the latent information flow, ``even in very large samples of 10,000 observations''. On 1997 NYSE TAQ data for the same two stocks, Cisco and Intel (Ané and Geman's own data were not available to them), and on other stocks, returns conditioned on trades or volume are ``not approximately Gaussian'' \citep[version of December 2006, abstract and secs.~1 and~4--6]{murphy-izzeldin-2006}. The trade-clock normality claim is therefore contested on the same stocks and year. Neither statement is about independence: ``IID'' is López de Prado's gloss.

His own support for dollar bars is one series: on E-mini S\&P~500 futures, the exponentially weighted number of tick and volume bars per day for a fixed bar size ``will vary wildly over the years'', and that of dollar bars less (Figure~2.1, p.~28). His exercises ask the reader to compare bar types by weekly bar-count stability, serial correlation, variance of monthly variances and Jarque--Bera statistics (Exercise~2.1, pp.~40--41): these are the tests to run on one's own data.

\qffield{Pitfalls} Better distributional properties are an empirical, market- and period-specific finding, not a property guaranteed by the sampling rule. Robustness of dollar bars to splits and corporate actions requires qualification; López de Prado assumes splits are already adjusted and says only that dollar bars ``tend to be robust'' to issuance and buybacks (p.~28). The ``Sampling as a function of trading activity creates returns closer to IID Normal'' attached López de Prado's gloss to Ané and Geman without the qualification above. Interpretation: an activity clock moves the volatility clustering into the bar durations, so a bar series has irregular calendar timing that must be kept for any time-of-day feature, for joins with other data and for the timing of decisions.

\qffield{See also} \hyperref[meth:information-driven-bars]{information-driven bars}; \hyperref[meth:event-based-sampling]{event-based sampling}; \hyperref[data:derived-market-statistics]{derived market statistics} (Velu's binned data); \hyperref[def:volatility-components]{volatility components}; \hyperref[ev:intraday-impact-profile]{intraday impact profile}.

\end{qfmethod}

\begin{qfmethod}[Information-driven bars: the tick rule, tick, volume and dollar imbalance bars, and tick, volume and dollar runs bars, with causal thresholds]\label{meth:information-driven-bars}
\qffield{Statement} Information-driven bars ``sample more frequently when new information arrives to the market'', information being meant in the microstructure sense of persistent imbalances in signed volume, associated with informed traders \citep[chap.~2, sec.~2.3.2, p.~29]{alpha-afml}. For ticks $(p_t,v_t)$ the \emph{tick rule} sets $b_t=b_{t-1}$ if $\Delta p_t=0$ and $b_t=\Delta p_t/|\Delta p_t|$ otherwise, so $b_t\in\{-1,1\}$, with $b_0$ ``set to match the terminal value $b_T$ from the immediately preceding bar'' (p.~29). Each bar closes at the first $T$ at which a statistic $\theta_T$ of the ticks since the bar opened reaches its expectation as set at the bar's start (pp.~29--32):
\begin{center}\footnotesize
\begin{tabular}{@{}lll@{}}
\toprule
Bar & $\theta_T$ & bar closes at the first $T$ with \\
\midrule
tick imbalance (TIB) & $\sum_{t=1}^Tb_t$ & $|\theta_T|\ge E_0[T]\,|2P[b_t=1]-1|$ \\
volume or dollar imbalance (VIB, DIB) & $\sum_{t=1}^Tb_tv_t$ & $|\theta_T|\ge E_0[T]\,|2v^+-E_0[v_t]|$ \\
tick runs (TRB) & $\max\{\sum_{t|b_t=1}b_t,\ -\sum_{t|b_t=-1}b_t\}$ & $\theta_T\ge E_0[T]\max\{P[b_t=1],1-P[b_t=1]\}$ \\
volume or dollar runs (VRB, DRB) & $\max\{\sum_{t|b_t=1}b_tv_t,\ -\sum_{t|b_t=-1}b_tv_t\}$ & $\theta_T\ge E_0[T]\max\{v^+,\ v^-\}$ \\
\bottomrule
\end{tabular}
\end{center}
with $v_t$ the units (VIB, VRB) or the currency amount (DIB, DRB) traded, $v^+=P[b_t=1]E_0[v_t|b_t=1]$ and $v^-=P[b_t=-1]E_0[v_t|b_t=-1]$, so that $E_0[v_t]=v^++v^-$ and $E_0[\theta_T]=E_0[T](v^+-v^-)=E_0[T](2v^+-E_0[v_t])$ for the volume imbalance (p.~30).

\qffield{Notation} $E_0[T]$ is the expected number of ticks per bar, $P[b_t=1]$ the unconditional probability that a tick is a buy. Imbalance bars offset buys against sells; runs bars count each side ``without offsetting them'' and allow sequence breaks, so a run is the larger one-sided total in the bar, not the longest unbroken sequence (p.~31).

\qffield{Conditions} The expectations are estimated as exponentially weighted moving averages ``from prior bars'': of $T$, of $b_t$ (for $2P[b_t=1]-1$), of $b_tv_t$, of the proportion of buy ticks and of buy and sell volumes (pp.~29--32). The threshold is therefore fixed when a bar opens and uses only closed bars, so bar formation is causal. López de Prado gives no EWMA span, no warm-up rule for the first bar and no $b_0$ for the very first tick: these are implementation choices to record.

Interpretation, derived here: the factorisation $E_0[\sum_{t\le T}x_t]=E_0[T]\,E_0[x_t]$ needs $T$ independent of the tick signs and sizes, or a Wald-type identity for IID ticks, neither of which the text states; and for runs bars $E_0[\max\{X,Y\}]\ge\max\{E_0X,E_0Y\}$, so the runs ``expectation'' is a heuristic threshold, not the expectation of $\theta_T$.

\qffield{Practice} López de Prado's interpretation is that TIBs ``are produced more frequently under the presence of informed trading'' and can be understood ``as buckets of trades containing equal amounts of information'' (p.~30); VIBs and DIBs, like volume and dollar bars, address tick fragmentation and outliers and, because the bar size adjusts, corporate actions (p.~31); runs bars target large traders who sweep the book, use icebergs or slice orders (p.~31). No test of these claims is given. In a toy order flow computed for this entry (\nolinkurl{afml-checks.py}), all four bar types closed more often during a one-sided burst than elsewhere, but on a perfectly balanced flow the tick-imbalance threshold $E_0[T]|2P[b_t=1]-1|$ is near zero, bars close after a tick or two, and the EWMA of $T$ then shrinks further: the scheme needs a floor or cap that the text does not give.

\qffield{Pitfalls} Putting a minus sign inside a one-element maximum in the runs statistic incorrectly makes the tick-runs $\theta_T$ equal to $T$ and the volume-runs $\theta_T$ the total volume; the source has two arguments, buy side against sell side (pp.~31--32). They also dropped the absolute values in the imbalance thresholds, so that ``$E_0[T]\,\vert\,2P[b_t=1]-1$'' read as a conditional expectation, and used $v^+$ without defining it. The tick rule classifies a tick by the last price change, not by the aggressor; where aggressor flags exist they are the better sign.

\qffield{See also} \hyperref[meth:activity-sampled-bars]{activity-sampled bars}; \hyperref[meth:event-based-sampling]{event-based sampling}; \hyperref[def:informed-trader-types]{informed trader types}; \hyperref[def:order-exposure-decisions]{order exposure decisions} (icebergs and slicing).

\qffield{Further reading} López de Prado's chap.~19 on microstructural features, not read here.

\end{qfmethod}

\begin{qfmethod}[The ETF trick: a basket, spread or rolled futures position as the value of \$1 invested, with holdings, costs, tradable volume and the single-futures roll]\label{meth:etf-trick}
\qffield{Statement} The ETF trick models ``a basket of securities as if it was a single cash product'', transforming a multi-product dataset ``into a single dataset that resembles a total-return ETF'', so that a strategy can always assume it trades ``cashlike products (non-expiring cash instruments)'' \citep[chap.~2, secs.~2.4--2.4.1, pp.~32--33]{alpha-afml}. It addresses spreads whose weights change (the spread can converge with no price change), spreads that go negative, constituents that do not trade at the same times, and execution costs (p.~33). With raw opens $o_{i,t}$ and closes $p_{i,t}$, the USD value of one point $\varphi_{i,t}$ (including FX), volume $v_{i,t}$ and carry, dividend or coupon $d_{i,t}$ (which can also charge margin or funding costs), allocations $\omega_t$ rebalanced or rolled on bars $B$, and $K_0=1$ (p.~34):
\[ h_{i,t}=\begin{cases}\dfrac{\omega_{i,t}K_t}{o_{i,t+1}\varphi_{i,t}\sum_i|\omega_{i,t}|}&t\in B\\ h_{i,t-1}&\text{otherwise}\end{cases},\qquad \delta_{i,t}=\begin{cases}p_{i,t}-o_{i,t}&(t-1)\in B\\ \Delta p_{i,t}&\text{otherwise,}\end{cases} \]
\[ K_t=K_{t-1}+\sum_ih_{i,t-1}\varphi_{i,t}(\delta_{i,t}+d_{i,t}). \]
Every instrument must have been tradable at $t-1$ and $t$ (p.~33). With $\tau_i$ the cost of trading \$1 of instrument $i$, the rebalance cost $c_t=\sum_i(|h_{i,t-1}|p_{i,t}+|h_{i,t}|o_{i,t+1})\varphi_{i,t}\tau_i$ for $t\in B$ is kept out of $K_t$ and treated as a negative dividend; the cost of buying or selling one unit is $\tilde c_t=\sum_i|h_{i,t-1}|p_{i,t}\varphi_{i,t}\tau_i$; and the tradable basket volume is $v_t=\min_i v_{i,t}/|h_{i,t-1}|$ (p.~34).

\qffield{Notation} $h_{i,t}$ holdings in securities or contracts; $\delta_{i,t}$ the change in market value per unit from $t-1$ to $t$. Dividing by $\sum_i|\omega_{i,t}|$ ``de-levers'' the allocation. Interpretation, derived: $\sum_i|h_{i,t}|o_{i,t+1}\varphi_{i,t}=|K_t|\sum_i|\omega_{i,t}|\operatorname{sgn}(o_{i,t+1}\varphi_{i,t})/\sum_j|\omega_{j,t}|$, which is $|K_t|$ when every $o_{i,t+1}\varphi_{i,t}>0$, and $K_t$ itself only when also $K_t\ge0$. With a negative price or conversion factor the sign of a holding no longer follows the sign of its weight and the sum is not a gross exposure.

\qffield{Conditions} \emph{Timing.} Holdings decided at $t$ are sized at the next bar's open $o_{i,t+1}$ because at a roll ``we may not know $p_{i,t}$ of the new contract'' (p.~34). Interpretation: this is an accounting convention, not an order. The quantity $h_{i,t}$ depends on $o_{i,t+1}$, which is unknown when the decision is taken at $t$, so weights $\omega_t$ that use only information up to $t$ are necessary but not sufficient for the series to be executable: an order placed at $t$ cannot carry that quantity.

The convention is reproduced only by an order whose quantity is set by its own fill price, for example a currency-amount (notional) order for $\omega_{i,t}K_t/\sum_i|\omega_{i,t}|$ in each leg, filled in full in a single-price opening auction at exactly $o_{i,t+1}$ and with $\varphi_{i,t}$ known (or the FX conversion fixed at that time); whether a venue accepts such orders, and fills them without rounding to lots, partial fills or a price other than the published open, must be checked for each instrument.

A quantity sized at $t$ on $p_{i,t}$ and filled at the open leaves a gross exposure different from $K_t$ when the open gaps (checked on a two-leg toy with opens 1.5 percent above and 1.6 percent below the closes: 1.0026 against $K_t=1$), and a trader who first observes the opening print can no longer trade at it. The gap between $p_{i,t}$ and $o_{i,t+1}$ on the old holdings is not booked, which is the intent at a roll (two contracts) and an approximation at a same-instrument rebalance, and $\varphi_{i,t}$ is applied to a $t+1$ price.

\emph{Positivity and capital.} López de Prado says the series will be ``strictly positive (at worst, infinitesimal)'' because profits and losses are reinvested at rebalances (pp.~33--34). Interpretation, checked: that holds for long-only baskets with positive prices, but a short leg can lose more than $K$ between rebalances (a short worth \$1 that rises 120 percent leaves $K=-0.2$). The recursion has no rule for exhausted capital: at $K_t=0$ the next holdings are zero, and at $K_t<0$ the formula for $h_{i,t}$ reverses every position (the long leg becomes short) while their gross value is $|K_t|$. A series that reaches $K_t\le0$ therefore stops describing the strategy; a backtest must stop it or state an explicit financing, margin or loss rule, which the method does not supply.

\emph{Costs.} $c_t$ charges the full value of old and new holdings, so it overstates the cost when holdings change little; López de Prado notes that cost functions need not be linear and can be simulated (p.~34). $v_t$ is undefined for a zero holding, which must be skipped.

\qffield{Practice} \emph{Single futures roll.} For one contract the ETF trick is a one-legged spread; López de Prado's more direct method forms the cumulative series of roll gaps, each the new contract's open minus the old contract's previous close, and subtracts it from the price series, backward (the end matches the raw series) or forward (the start matches) (sec.~2.4.3, pp.~36--37, Snippet~2.2). Rolled prices serve for PnL and mark-to-market; ``raw prices should still be used to size positions and determine capital consumption'', and rolled prices ``can indeed become negative, particularly in futures contracts that sold off while in contango'' (p.~37). The non-negative alternative computes the return as the rolled price change divided by the \emph{previous raw price} and compounds $(1+r)$ (p.~37, Snippet~2.3).

Interpretation: the result stays non-negative while $1+r\ge0$ on every bar, which holds on non-roll bars when raw prices are positive but not if raw prices themselves turn negative. A two-leg toy (\nolinkurl{afml-checks.py}) is a consistency check only: the change in $K$ equals the holdings' profit and loss recomputed from the same prices, and gross exposure at the booked open equals $K$ (here $K>0$ and all prices positive). The timing and capital cases above are computed in \nolinkurl{financial-labels-capital-checks.py}.

\qffield{Pitfalls} Use the next opening price $o_{i,t+1}$ in $h_{i,t}$ and the previous raw price in the return denominator. Positivity requires the capital and exposure conditions stated above. The allocation vector itself is a separate choice; one way to set it is in \hyperref[mod:pca-risk-distribution-weights]{PCA risk-distribution weights}.

\qffield{See also} \hyperref[mod:pca-risk-distribution-weights]{PCA risk-distribution weights}; \hyperref[conv:cme-globex-trading]{CME Globex trading} (contract months); \hyperref[def:futures-hedge-basis]{futures hedges and basis} (stack and roll); \hyperref[def:transaction-costs]{transaction costs}; \hyperref[mech:etf-creation-redemption]{ETF creation and redemption} (real ETFs, as against this virtual one).

\end{qfmethod}

\begin{qfmethod}[Event-based sampling: downsampling for reduction and the symmetric CUSUM filter, with its update, reset and threshold rules]\label{meth:event-based-sampling}
\qffield{Statement} Fitting on every bar is usually a poor idea, because some ML algorithms do not scale with sample size and because they ``achieve highest accuracy when they attempt to learn from relevant examples'' \citep[chap.~2, sec.~2.5, p.~38]{alpha-afml}. \emph{Sampling for reduction} (downsampling) takes bars at a constant step (linspace sampling), which is simple but has an arbitrary step and depends on the seed bar, or uniformly at random; neither selects the most informative observations (sec.~2.5.1, p.~38). \emph{Event-based sampling} keeps the bars at which some event occurs and lets the model learn whether a prediction function exists there; if not, the event is redefined (sec.~2.5.2, p.~38). The \emph{CUSUM filter}, a quality-control method for detecting a shift in a mean, sets for observations $y_t$
\[ S_t=\max\{0,\,S_{t-1}+y_t-E_{t-1}[y_t]\},\quad S_0=0, \]
and recommends an action at the first $t$ with $S_t\ge h$, the filter size; the symmetric filter adds $S_t^-=\min\{0,S^-_{t-1}+y_t-E_{t-1}[y_t]\}$ and uses $S_t=\max\{S_t^+,-S_t^-\}$ (sec.~2.5.2.1, pp.~38--39). For sampling, ``we will sample a bar $t$ if and only if $S_t\ge h$, at which point $S_t$ is reset'' (p.~39).

\qffield{Conditions} $S_t=0$ whenever $y_t\le E_{t-1}[y_t]-S_{t-1}$: the zero floor skips downward deviations, so the one-sided filter detects run-ups from the last reset (p.~39). Without a reset, $S_t\ge h$ if and only if some partial sum $\sum_{i=\tau}^t(y_i-E_{i-1}[y_i])\ge h$ (López de Prado prints $E_{i-1}[y_t]$); interpretation, after a reset $\tau$ must start after it. López de Prado describes the inputs as ``IID observations arising from a locally stationary process'' (pp.~38--39), but his Snippet~2.4 sets $E_{t-1}[y_t]=y_{t-1}$ and cumulates first differences of a level series such as a price, which treats $y$ as a level with martingale-like increments; the operational definition is the cumulated deviation $y_t-E_{t-1}[y_t]$ under whatever forecast is used. In the snippet the comparisons are strict ($S^-<-h$, $S^+>h$), only the side that triggered is reset, and at most one event fires per bar, the negative side checked first (p.~39).

\qffield{Practice} A filter event needs ``a full run of length $h$'', so a series hovering near a level does not fire repeatedly, unlike Bollinger bands; $S_t$ can be built on other features, such as structural-break statistics or microstructure measures (p.~40). Checked on toy series (\nolinkurl{afml-checks.py}): a series hovering within $\pm0.05$ of a level fires no event at $h=0.5$, and a run of exactly $h$ fires under the text's $\ge$ but not under the snippet's $>$. López de Prado's exercises use returns, or absolute returns, as $y_t$ with $h=0.05$ (p.~40); interpretation, with the snippet's differencing that needs a cumulative return series as input. Lam and Yam's CUSUM trading rule and Fama and Blume's filter rule are different uses: they trade on the signal, whereas here the filter only selects observations (p.~39).

\qffield{Pitfalls} The zero condition uses $\le$, not equality alone. The run condition must use the correct expectation index and apply the sampling and reset rule. The threshold $h$ and any volatility scaling of it (López de Prado's labelling exercise sets $h$ to the standard deviation of daily returns, p.~55) are parameters: estimating them on the full sample before a backtest uses future data; see \hyperref[pit:label-weight-leakage]{time-validity of labels and weights}.

\qffield{See also} \hyperref[meth:triple-barrier-labelling]{triple-barrier labelling} (which the sampled events seed); \hyperref[meth:activity-sampled-bars]{activity-sampled bars}; \hyperref[meth:information-driven-bars]{information-driven bars}.

\qffield{Further reading} Structural-break and entropy features (López de Prado, chaps.~17--18) are not covered.

\end{qfmethod}
